% SPDX-License-Identifier: CC-BY-4.0
\documentclass[11pt,letterpaper]{article}

\usepackage[T1]{fontenc}
\usepackage{amsmath,amssymb}
\usepackage{tikz}
\usepackage{pgfplots}
\pgfplotsset{compat=1.18}
\usepackage{algorithm}
\usepackage{algpseudocode}
\usepackage{booktabs}
\usepackage{enumitem}
\usepackage[hidelinks]{hyperref}

\newcommand{\N}{\mathbb{N}}
\newcommand{\Z}{\mathbb{Z}}
\newcommand{\Q}{\mathbb{Q}}
\newcommand{\R}{\mathbb{R}}
\newcommand{\C}{\mathbb{C}}
\newcommand{\Map}[3]{#1 \colon #2 \to #3}
\newcommand{\Mod}[1]{#1\text{-}\mathrm{Mod}}
\newcommand{\GroupHom}{\mathrm{Hom}_{\mathbf{Grp}}}
\newcommand{\RingHom}{\mathrm{Hom}_{\mathbf{Ring}}}

\newtheorem{theorem}{Theorem}
\newtheorem{lemma}[theorem]{Lemma}
\newtheorem{corollary}[theorem]{Corollary}
\newtheorem{definition}[theorem]{Definition}
\newtheorem{remark}[theorem]{Remark}
\newenvironment{proof}{\noindent\textit{Proof.}}{\hfill\(\square\)}

\title{Congruence transport of affine states}
\author{ascendnoosphere}
\date{\today}

\begin{document}
\maketitle

\begin{abstract}
Let \(T_j(q)\) be the sum of the powers \(q^0,\dots,q^{j-1}\), with the empty
sum equal to zero, and let
\(A_j(q,c,n)=q^j n+c\,T_j(q)\). Replacing the seed \(n\) by \(n+td\) adds
exactly \(q^j td\) to this state, for every pair of integers \(t,d\) and every
time \(j\). The difference is therefore a multiple of \(d\) at every time, and
divisibility by \(d\) is unchanged at every time. When \(d\) is the signed
divisor of a parameter record, the first time at which the divisor divides the
affine state is the same for \(n\) and for \(n+td\). The argument is an
identity in \(\Z\). It does not produce a state of large absolute value.
\end{abstract}

\section{Affine states}

Throughout, \(q,c,n,t,d\) denote integers and \(j\) denotes a natural number,
including zero. The offset sum is defined by \(T_0(q)=0\) and
\[
T_{j+1}(q)=T_j(q)+q^j.
\]
Thus \(T_j(q)=q^0+\cdots+q^{j-1}\) when \(j\ge 1\). No closed form of this sum
is required below: the sum depends on \(q\) and \(j\), not on the seed.

\begin{definition}
\label{def:affine-state}
The \emph{affine state} at time \(j\) is
\[
A_j(q,c,n)=q^j n+c\,T_j(q).
\]
\end{definition}

A parameter record consists of integers \(q\), \(c\), and \(d\) with
\(|d|\ge 2\). In particular \(d\neq 0\), since \(|d|\ge 2>0\). The raw map
sends \(n\) to \(n/d\) when \(d\) divides \(n\), and to \(qn+c\) otherwise.
The theorems in this note do not iterate that map. They concern only the
integer \(A_j(q,c,n)\) and divisibility of that integer.

\begin{definition}
\label{def:first-hit}
Let \(p\) be a parameter record with multiplier \(q\), offset \(c\), and
divisor \(d\). The time \(j\) is a \emph{first hit} of the seed \(n\), written
\(\operatorname{Hit}(p,n,j)\), when
\[
d\mid A_j(q,c,n)
\]
and
\[
d\nmid A_i(q,c,n)\qquad\text{for every natural number \(i<j\).}
\]
\end{definition}

When \(j=0\), the second requirement is vacuous, so a first hit at time zero
is exactly divisibility of the seed. The predicate does not assert that a raw
orbit equals the affine state, and it does not assert that any seed has a
first hit.

\section{The shift identity}

\begin{theorem}
\label{thm:affine-class-shift}
For all \(q,c,n,t,d\in\Z\) and every \(j\in\N\),
\[
A_j(q,c,n+td)=A_j(q,c,n)+q^j td.
\]
\end{theorem}

\begin{proof}
Expand the left-hand side by Definition~\ref{def:affine-state}:
\[
A_j(q,c,n+td)=q^j(n+td)+c\,T_j(q).
\]
Distributivity in \(\Z\) gives \(q^j(n+td)=q^j n+q^j(td)\), and multiplication
of integers is associative, so \(q^j(td)=(q^j t)d=q^j td\). Therefore
\begin{align*}
A_j(q,c,n+td)
&=q^j n+q^j td+c\,T_j(q)\\
&=(q^j n+c\,T_j(q))+q^j td\\
&=A_j(q,c,n)+q^j td,
\end{align*}
where the last step is again Definition~\ref{def:affine-state}. The offset
sum is the same on both sides because it does not depend on the seed.
\end{proof}

The identity is an equality of integers, not merely a congruence. It holds
for every integer \(d\), including zero and the units \(\pm 1\), and for
negative \(q\), \(t\), and \(d\). Admissibility \(|d|\ge 2\) is not used.

\begin{theorem}
\label{thm:difference-multiple}
For all \(q,c,n,t,d\in\Z\) and every \(j\in\N\),
\[
d\mid A_j(q,c,n+td)-A_j(q,c,n).
\]
In particular, if \(i<j\), the same statement at time \(i\) gives
\[
d\mid A_i(q,c,n+td)-A_i(q,c,n).
\]
\end{theorem}

\begin{proof}
Theorem~\ref{thm:affine-class-shift} at time \(j\) rearranges to
\[
A_j(q,c,n+td)-A_j(q,c,n)=q^j td.
\]
The right-hand side equals \(d\cdot(q^j t)\), so \(d\) divides the difference.
The natural number \(j\) was arbitrary. Applying the same conclusion at an
earlier time \(i\) gives the second claim.
\end{proof}

\begin{theorem}
\label{thm:divisibility-shift}
Let \(p\) be a parameter record with multiplier \(q\), offset \(c\), and
divisor \(d\). For every \(n,t\in\Z\) and every \(i\in\N\),
\[
d\mid A_i(q,c,n)
\quad\text{if and only if}\quad
d\mid A_i(q,c,n+td).
\]
\end{theorem}

\begin{proof}
Write \(S=A_i(q,c,n+td)\) and \(U=A_i(q,c,n)\). By
Theorem~\ref{thm:difference-multiple} there is an integer \(k\) such that
\[
S-U=d k.
\]
Both directions are witnessed explicitly.

Suppose first that \(d\) divides \(U\). Then \(U=dm\) for some \(m\in\Z\).
Substituting the difference gives
\[
S=U+(S-U)=dm+dk=d(m+k),
\]
so \(d\) divides \(S\).

Suppose conversely that \(d\) divides \(S\). Then \(S=dm\) for some
\(m\in\Z\). Solving the difference for the unshifted state gives
\[
U=S-(S-U)=dm-dk=d(m-k),
\]
so \(d\) divides \(U\).

The two constructions are inverses as integer witnesses, but the proof uses
only that each direction produces some witness. The time \(i\) is arbitrary,
so the equivalence holds at the hitting time and at every earlier time.
\end{proof}

\section{Transport of the first hit}

\begin{theorem}
\label{thm:first-hit-shift}
Let \(p\) be a parameter record with divisor \(d\). For every \(n,t\in\Z\) and
every \(j\in\N\),
\[
\operatorname{Hit}(p,n,j)
\quad\text{if and only if}\quad
\operatorname{Hit}(p,n+td,j).
\]
\end{theorem}

\begin{proof}
Let \(q\) and \(c\) be the multiplier and offset of \(p\). By
Definition~\ref{def:first-hit}, the forward implication assumes
\[
d\mid A_j(q,c,n)
\]
and \(d\nmid A_i(q,c,n)\) for every \(i<j\), and must produce the same two
statements with seed \(n+td\). The backward implication assumes the two
statements for \(n+td\) and must return them for \(n\).

Apply Theorem~\ref{thm:divisibility-shift} at time \(j\). Divisibility of
\(A_j(q,c,n)\) is equivalent to divisibility of \(A_j(q,c,n+td)\), so the
hitting clause transports in both directions.

It remains to transport the failures, also in both directions. Fix \(i<j\).
Theorem~\ref{thm:difference-multiple} supplies a multiple of \(d\) at this
earlier time, and Theorem~\ref{thm:divisibility-shift} turns that multiple
into the equivalence
\[
d\mid A_i(q,c,n)
\quad\text{if and only if}\quad
d\mid A_i(q,c,n+td).
\]
Thus failure for \(n\) implies failure for \(n+td\), and failure for
\(n+td\) implies failure for \(n\). Since \(i<j\) was arbitrary, the universal
quantifier over earlier times is preserved. Both directions of
Definition~\ref{def:first-hit} follow.
\end{proof}

\begin{remark}
\label{rem:not-a-large-state}
The equivalence identifies two divisibility patterns. It does not assert that
either pattern occurs, and the correction term \(q^j td\) may be zero, for
instance when \(t=0\) or \(j>0\) and \(q=0\). Nothing in the identity forces
\(|A_j(q,c,n)|\) to exceed a prescribed bound. Producing a large state, if it
is possible, requires a separate choice of seed or shift and a separate
estimate.
\end{remark}

\begin{remark}
\label{rem:not-the-raw-orbit}
If some earlier raw iterate is divisible by \(d\), the raw map leaves the
affine branch, while \(A_i(q,c,n)\) is still defined. Equality of the raw
iterate with the affine state needs a division-free prefix and is not proved
here. The transport applies to the affine state whether or not that prefix
holds.
\end{remark}

\end{document}
